\chapter{Progettazione}
\label{cap:progettazione}
In questo capitolo verrà trattato il secondo periodo del tirocinio: la progettazione e la codifica. Questa parte, avvenuta dopo l'Analisi dei Requisiti, è da ritenersi importante per il progetto in quanto ha occupato gran parte del tempo lavorativo. Il primo passo è stato eseguire la progettazione dell'architettura, in modo che i servizi necessari fossero correttamente configurati prima del loro effettivo utilizzo. Successivamente è stata eseguita la stesura del codice per la realizzazione della Skill. Nei paragrafi successivi seguiranno porzioni di codice significativo per la loro importanza e per il funzionamento di determinati servizi. Importante ricordare che il codice implementato e mostrato sarà in Node.js come quanto detto nel paragrafo \hyperref[nodejs]{1.4.1}. 

\section{Architettura}
Come riportato nel paragrafo \hyperref[serivizi-aws]{1.4.3}, il progetto si basa interamente sui servizi offerti dall'ecosistema Amazon. Di conseguenza è risultato facile ed immediato impostare i servizi dell'architettura in modo da rendere il tutto efficacie ed efficiente. In questo paragrafo si andrà a spiegare cosa offre ogni singolo servizio di AWS e come esso è stato impostato per il suo corretto funzionamento nella Skill.
